% Source: https://artofproblemsolving.com/community/c3h2495432_regular_heptagon
% Stable identity: aops-topic-2495432-q01
Let $ABCDEFG$ be a regular heptagon, and let $x=\angle ADC$ and $y=\angle DAB$.
\begin{enumerate}
\item[(a)] Prove that $x=y$, find this common value, and deduce that $ABCD$ is an isosceles trapezoid with $AD\parallel BC$.
\item[(b)] A student tries to prove $x=y$ as follows. Each interior angle of the heptagon is $\dfrac{900^\circ}{7}$. Triangle $ADE$ is isosceles with $AD=AE$, and its base angles are $\angle ADE=\angle AED=\dfrac{900^\circ}{7}-x$, so its vertex angle is $b=\angle EAD=180^\circ-2\left(\dfrac{900^\circ}{7}-x\right)$. By symmetry $\angle EAG=\angle DAB=y$, so $2y+b=\dfrac{900^\circ}{7}$. But substituting $b$ into this equation does not give $x=y$. Explain why this approach cannot work.
\end{enumerate}
